\section{Scope, conventions, and present conclusions}

The principal target is $K_{4k}(\mathbf Z)=0$ for every positive
integer $k$, and hence Vandiver's conjecture at every prime.
That target remains open in this research snapshot.
The checked results include new complete finite certificates for
$K_{52}(\mathbf Z)=K_{56}(\mathbf Z)=K_{60}(\mathbf Z)=K_{64}(\mathbf Z)=K_{68}(\mathbf Z)=K_{72}(\mathbf Z)=K_{76}(\mathbf Z)=K_{80}(\mathbf Z)=0$, uniform dimension and
spectrum formulas, an independent computation of $157899$, a closed
source extraction of the archive's nine-hypothesis implication, and
\[
 \boxed{\ \log|K_{4k}(\mathbf Z)|<\frac14(2k)^4\quad(k\ge1).\ }
\]
This order bound is not a vanishing theorem. The completed finite
degrees, together with K17, give $K_{4k}(\mathbf Z)=0$ for
$1\le k\le20$. The all-weight arithmetic elimination remains
active. All finite computations started in this continuation are complete;
no higher degree is claimed.

The task's source convention is retained throughout: Alp\"oge's
\emph{Remark}, step~17, and the pinned Lean archive are cited as
theorems. This report does not constitute a new validation of the
archive's entire foundational or imported theorem chain. The fresh
source extraction is distinguished from a Lean kernel run: no Lean
executable was available and no build or new axioms report was run.
The source files actually retrieved have pinned Git-blob provenance.

\paragraph{Carried results.}
The supplied K17 paper and the intake reports settle the following
inputs, which are cited rather than re-derived:
K17-S17-UNIT; the arithmetic comparison VK-PT, (K1)--(K5), and
VK-EXCEPTIONS; K17-U3-HECKE's lattice class number and integral
harmonic model; K17-GL2-GENERAL's integral level-three basis;
K17-MINUS7-GENERAL-SOURCE; K17-U3-RESIDUAL-RING; and the finite
vanishings through $K_{48}$. Exact arithmetic comparison references
are Weibel, \emph{The K-book}, VI.8.2, VI.8.3, VI.10.9, and
Kurihara, \emph{Some remarks on conjectures about cyclotomic fields
and $K$-groups of $\mathbf Z$} (1992), Lemma~1.2, Corollary~1.5,
Proposition~3.7 and Remark~3.11, as reconciled in K17.
The modern comparison is required for the two-way zero-equivalence;
no unconditional 1992 proof of the then-conjectural comparison is
being attributed to Kurihara.

At odd primes below $2^{31}$ we use Hart--Harvey--Ong,
\emph{Irregular primes to two billion}, \S1, item~(2), as carried
by K17. Prime two is separate: Rognes--Weibel,
\emph{Two-primary algebraic $K$-theory of rings of integers in
number fields}, Theorem~0.7, proved as Theorem~7.11,
with K17's localization specialization, gives the required
$2$-primary vanishing. The large-modulus calculations never
assume a residual integer is prime.

For the geometric questions AN\_GRH A05, A33 and A46, and the
previous failed route-two constructions A36--A44, are carried at
their stated intake scope. K17-VANDIVER-KERNEL,
K17-GREENBERG-KERNEL, K17-P2-MOMENT, and K17-CUBIC-ODD retain
exactly their stated arithmetic hypotheses. No rational
mixed-Tate statement is substituted for an integral hypothesis.

\paragraph{Reading the paper and evidence.}
Part~I contains checked findings, including negative answers and
source comparisons. Part~II records unresolved questions and
proposals. A theorem that invokes a named background theorem is
checked as an implication from that background, rather than by
re-proving the entire literature. The numerical evidence includes
complete matrices, exact finite arithmetic, input hashes, and
scripts. Producer completion, a full independent replay, a
mathematical proof reading, and a formal kernel check are different
events in \path{CHECK_EVENTS.tsv}.

All file locators within the supplied K17, AN\_GRH and GRH\_NIGHT
returns refer to the exact attached versions. Scan numbers refer
to cote~149's indexing as preserved in those supplied transcripts;
no fresh palaeographic reading of every original sheet is claimed.
The registers \path{QUESTIONS.tsv}, \path{LEDGER.md}, and
\path{IMPLICATIONS.md} state the live boundaries of each conclusion.
